% Introduction (PUB-REF-0001 AN-21 F-BOR: what the product is and does,
% the manuals and their jobs with routing, installing, typefaces, contact).
% Sources (Law 1): ADR-0007 DEC-01/02/05/08; os/tps/tps.pas (TPSIN.PAS,
% TPSOUT.S, output on the serial port through the RTL); seed/rt/start_dos.asm
% (COM1); docs/pubs/originals/borland-turbo-pascal.md s1, s3.
\Gates{test-pubs-build, test-tps-gen-os, test-compiler-os, test-pubs-tps}
\IntroductionChapter

Welcome to Turbo Initech\reg{}, the Pascal compiler for InitechOS.  Turbo
Initech is small, fast to learn and exact: it accepts a compact Pascal,
checks every program completely before it writes a single instruction,
and turns the program into 80386 code that runs under InitechDOS\tm{}.
Turbo Initech is written in its own language, and it compiles itself.

This book shows you how to use Turbo Initech.  Chapters 1 through 5 teach
you to write, compile, run and correct programs; Chapter 8 is the
reference for the compiler's command and for everything it prints.

\section{What is in Turbo Initech}
\begin{itemize}
\item The compiler \file{TPS.COM}, which reads a program from the file
  \file{TPSIN.PAS} and writes its translation to \file{TPSOUT.S}.
\item A language with integers, characters, Boolean values, strings,
  constants, arrays, records, procedures and functions with value and
  \rw{var} parameters, recursion, and files.
\item A run-time library that writes to the serial port and reads and
  writes files through InitechDOS.
\item Listings of every pass of the compiler: the tokens, the structure
  of the program, and the layout of every variable.
\end{itemize}

\section{What you need}
\ix{system requirements}
\begin{itemize}
\item A computer running InitechOS with InitechDOS, as described in the
  \doctitle{InitechDOS User's Guide}.
\item A terminal, or a second computer with a communications program,
  connected to the first serial port, COM1.  The compiler and the programs
  it compiles write their output to COM1, not to the display.
\item A text editor to write programs, such as the InitechDOS line editor
  \cmd{EDLIN}.
\end{itemize}

\section{The Turbo Initech manuals}
\ix{manuals}
\SideHead{User's Guide}This book (1077-0000-10) teaches the language by
example and describes the compiler.  Begin here.

\SideHead{Programmer's Guide}The \doctitle{Turbo Initech Programmer's
Guide} (1077-0000-11) is the definition of the language: its syntax, every
rule, the layout of data in memory, the code the compiler writes, and every
message.

\SideHead{Library Reference}The \doctitle{Turbo Initech Library Reference}
(1077-0000-12) describes each standard procedure and function.

\SideHead{Self-Compilation Guide}The \doctitle{Turbo Initech
Self-Compilation Guide} (1077-0000-13) describes how Turbo Initech is
built and how it proves that it reproduces itself.

\SideHead{Quick Reference}The \doctitle{Turbo Initech Quick Reference}
(1077-0000-20) summarizes the language on a few pages.

If you know Pascal already, read Chapter 1 of this book and then the
\doctitle{Programmer's Guide}, whose margin notes mark the places where
Turbo Initech differs from other Pascal compilers.  If you want to know
what a message means, look it up in the \doctitle{Programmer's Guide},
Appendix A.

\section{Installing Turbo Initech}
\ix{installing}
Turbo Initech needs no installation program: copy \file{TPS.COM} to the
directory in which you keep your programs.
\draft{the Turbo Initech distribution diskette is not produced, and
TPS.COM is not on the shipped InitechOS data diskette; it is built by the
genesis compiler (Self-Compilation Guide, Chapter 3)}

\section{Typefaces used in these books}
\ix{typefaces}
\begin{BorTable}{Typefaces}{p{38mm}p{62mm}}
\Hd{Example} & \Hd{Meaning}\EndHd
\texttt{Writeln(Total);} & Program text, file names, and what the compiler
  and your programs print\\
\texttt{\bfseries begin} & A reserved word inside program text\\
\textbf{begin} & A reserved word in running text\\
\textit{Count} & An identifier of your program, a placeholder, or a
  standard routine named in running text\\
\typed{TPS} & Something you type\\
\key{Enter} & A key on the keyboard\\
\MarginFace Margin note & A note on the paragraph beside it
\end{BorTable}

\section{How to contact Initech}
\ix{technical support}
If you have a question that the manuals do not answer, call the Initech
Help Desk at 555-0125, extension 2504, or write to \InitechAddressLine.
Please have the program, and the line the compiler printed, at hand.
